\subsection*{Partie A}

\begin{minipage}{0.7\linewidth}
	Soit $f$ la fonction définie pour tout nombre réel $x$ par $f(x) = (12x-480)(x-10)$.

	Établir le tableau de signes de la fonction $f$ sur $\R$.
\end{minipage}
\hfill
\fbox{
	\begin{minipage}{0.2\textwidth}
		\textbf{Aide au calcul :}\\[0.5em]
		$480 = 40 \times 12$
	\end{minipage}
}

\subsection*{Partie B}

On s'intéresse dans cette partie à une usine qui fabrique des tracteurs. Chaque mois elle en vend entre 0 et 18. On note $x$ le nombre de tracteurs vendus par mois.

On modélise alors le bénéfice de cette entreprise, \textbf{\uline{en centaines d'euros}}, par la fonction $B$ définie sur l'intervalle $[0~;~18]$ par 

\[ B(x) = 4x^3 - 300x^2 + \np{4800}x - \np{10000}. \]

\begin{enumerate}
	\item Calculer $B(0)$.
	\item On donne $B(18) = \np{2528}$. Interpréter ce résultat dans le contexte de l'exercice.
	\item
	\begin{enumerate}
		\item Pour tout réel $x$ de l'intervalle $[0 ~;~18]$, calculer $B'(x)$ puis montrer que $B'(x) = f(x)$, où $f$ est la fonction définie dans la partie A.
		\item En déduire le tableau de variations de la fonction $B$ sur l'intervalle $[0~;~18]$.
		\item À l'aide du tableau de variations, déterminer combien de tracteurs l'entreprise doit produire et vendre afin d'obtenir un bénéfice maximal. Quelle sera la valeur de ce bénéfice maximal ?
	\end{enumerate}
\end{enumerate}
